\documentclass[10pt,a4paper]{amsart}
\usepackage[margin=19mm]{geometry}
\usepackage{amsmath,amssymb,mathtools}
\usepackage[scr=rsfs,cal=euler]{mathalfa}
\usepackage{xfrac}
\usepackage[unicode,hidelinks,bookmarks=false]{hyperref}
\newcommand{\ie}{i.e.\ }
\newcommand{\cf}{cf.\ }
\newcommand{\resp}{respectively\ }
\newcommand{\Subsection}{\subsection}
\providecommand{\nobreakdash}{\nobreak\hbox{-}}
\setlength{\emergencystretch}{2em}
\multlinegap0pt
\usepackage{bm}
\usepackage{bbm}
\usepackage{dsfont}
\usepackage{xspace}
\usepackage{cancel}
\usepackage{mathtools}
\usepackage{amsthm}
\makeatletter
\@ifundefined{thm}{\newtheorem{thm}{Theorem}}{}
\makeatother
\title[Optimal K dual frames and pairs in the presence of erasures]{Optimal K dual frames and pairs in the presence of erasures}
\author{Shankhadeep Mondal, Deguang Han, R. N. Mohapatra}
\date{Source: arXiv:2508.01964v1 (CC BY 4.0)}
\begin{document}
\maketitle

\noindent\emph{Source and licence.} Optimal K dual frames and pairs in the presence of erasures. Version arXiv:2508.01964v1 (CC BY 4.0), distributed under the Creative Commons Attribution 4.0 International licence, as recorded in the arXiv licence metadata for this exact version and shown on the abstract page https://arxiv.org/abs/2508.01964v1. Full rights notice: licenses/arxiv-2508.01964-CC-BY-4.0.txt.\par\smallskip

\noindent\textit{Editorial scope.} Counted unit: the Thm whose source label is \texttt{theorem3point5twoerasure} in the article (source0.tex lines 336--355, source label \texttt{theorem3point5twoerasure}), delivered together with the article's own complete original proof of it (source0.tex lines 358--407, one \texttt{proof} environment of 50 lines). No mathematics is written, repaired or re-derived: statement and proof are the article's own text line for line. Required context retained uncounted: thm5point2 (lines 291--294); eqn3point4* (lines 329--331). Every definition, display and label of those lines is retained unchanged. Omitted from this package and not redistributed: 1--290, 295--328, 332--335, 356--357, 408--1476. Nothing inside a retained excerpt is omitted or altered: every retained body is delivered exactly as the article's lines are written, so no omission receipt of any kind accompanies this package. Citations: the retained excerpts carry no citation command at all, so no bibliography entry of the article is transcribed and no reference is redistributed. Reference adaptation: none was needed and none was made; every reference inside the delivered unit resolves inside it, so no unresolved reference or citation is produced. Printed slips kept literal: no printed word, symbol, hypothesis or number of the counted statement or of its proof was altered. Layout and class: the article's own class is \texttt{elsarticle} with packages including amssymb, amsmath, amsthm, inputenc, marginnote, amsfonts, textcomp, color, soul, placeins, tikz, makecell, graphicx, hyperref; the wrapper is the lane's pinned \texttt{amsart} editorial wrapper at 10pt on A4, which restates the theorem-like environments the retained text needs and adds the article's own macro definitions that the retained text needs (none), with extra wrapper packages (bm, bbm, dsfont, xspace, cancel, mathtools). The delivered numbering is the unit's own. Assets: no retained span contains a figure environment, a graphics-inclusion command, a PDF-page inclusion, a table or a TikZ picture; no image file is copied into this package, the package declares no asset dependency and no image file is redistributed with it. Licence: version arXiv:2508.01964v1 of this article is distributed under the Creative Commons Attribution 4.0 International licence, as recorded in the arXiv licence metadata for this exact version and shown on the abstract page https://arxiv.org/abs/2508.01964v1. A full rights notice is delivered at licenses/arxiv-2508.01964-CC-BY-4.0.txt. Characters: every retained excerpt is pure ASCII, so the delivered engine, XeLaTeX with its default Latin Modern text font, reports no missing character.
	\begin{thm} \label{thm5point2}
   Let $K \in \mathcal{B}(\mathcal{H}_n)$ be a positive semi-definite operator. Then, the value of $\tilde{r}_{1} $ is $\frac{tr(K)}{N}.$  Moreover, an $(N,n)$ dual pair $(F,G) \in \mathcal{R}_{1} $ if and only if it is a $1-$uniform $K-$dual pair.
		
	\end{thm}

\begin{align}\label{eqn3point4*}
	r_{2}(F,G) = \max_{i\neq j} \left| \dfrac{\operatorname{tr}(K)}{N} \pm \sqrt{\alpha_{ij}\alpha_{ji}} \right| = \max_{i\neq j} \left| \dfrac{\operatorname{tr}(K)}{N} + \sqrt{\alpha_{ij}\alpha_{ji}} \right|.
\end{align}

	\begin{thm}\label{theorem3point5twoerasure}
Let \( \mathcal{H}_n \) be a real Hilbert space of dimension \( n \) and let \( K \in \mathcal{B}(\mathcal{H}_n) \) be a positive semi-definite operator. Then the minimal spectral radius over all $K-$dual pairs for two erasures satisfies
\[
\tilde{r}_{2} \geq 
\begin{cases}
\dfrac{\operatorname{tr}(K)}{N} + \sqrt{\dfrac{N\,\operatorname{tr}(K^2) - \left(\operatorname{tr}(K)\right)^2}{N^2(N-1)}}, & \text{if } \operatorname{tr}(K^2) \geq \dfrac{\left(\operatorname{tr}(K)\right)^2}{N},\\~\\
\sqrt{\dfrac{(N-2)\left(\operatorname{tr}(K)\right)^2 + N\,\operatorname{tr}(K^2)}{N^2(N-1)}}, & \text{if } \operatorname{tr}(K^2) < \dfrac{\left(\operatorname{tr}(K)\right)^2}{N}.
\end{cases}
\]

Moreover, if there exists a  $2-$uniform \( K-\)dual pair, then the inequality above becomes an equality:
\begin{equation}\label{equation7expression}
\tilde{r}_{2} = 
\begin{cases}
\dfrac{\operatorname{tr}(K)}{N} + \sqrt{\dfrac{N\,\operatorname{tr}(K^2) - \left(\operatorname{tr}(K)\right)^2}{N^2(N-1)}}, & \text{if } \operatorname{tr}(K^2) \geq \dfrac{\left(\operatorname{tr}(K)\right)^2}{N},\\~\\
\sqrt{\dfrac{(N-2)\left(\operatorname{tr}(K)\right)^2 + N\,\operatorname{tr}(K^2)}{N^2(N-1)}}, & \text{if } \operatorname{tr}(K^2) < \dfrac{\left(\operatorname{tr}(K)\right)^2}{N}.
\end{cases}
\end{equation}
In this case, a dual pair \( (F, G) \in \mathcal{R}_2 \) if and only if it is a  $2-$uniform \( K \)-dual pair.
\end{thm}

	\begin{proof}
		Suppose \( (F, G) \) is an \( (N, n) \) dual pair. Taking \( \alpha_{ij} = \langle g_i, f_j \rangle \), it is easy to see that
\[
\sum_{i,j = 1}^N \alpha_{ij} \alpha_{ji} 
= \sum_{i=1}^N \left\langle g_i , \sum_{j=1}^N \langle f_i , g_j \rangle f_j \right\rangle 
= \sum_{i=1}^N \langle g_i , K f_i \rangle 
= \operatorname{tr}(\Theta_G^* \Theta_{K F}) 
= \operatorname{tr}\left({K^*}^2\right) 
= \operatorname{tr}(K^2).
\]
 Further, if $(F,G)$ is $1-$uniform $K-$dual pair, then
		\begin{align}\label{eqation7}
			\displaystyle{\sum_{i \neq j}\alpha_{ij}\alpha_{ji}} = tr\left({K^*}^2  \right)-\sum\limits_{i=1}^N \alpha_{ii}^2 = tr\left({K}^2  \right) - \frac{(tr(K))^2}{N}.
		\end{align}	
		Let  $\mu :=  tr\left({K}^2  \right) - \dfrac{(tr(K))^2}{N}.$ In the following, we derive a lower bound for \( \tilde{r}_{2} \), which varies according to the sign of \( \mu \).

		\par Suppose $\mu \geq 0.$ 	For $(F,G)\in \mathcal{R}_{1}, $   as in (\ref{eqn3point4*}), we have	$r_{2} (F,G) = \max\limits_{i \neq j} \left| \frac{tr(K)}{N}+ \sqrt{ \alpha_{ij}\alpha_{ji}} \right|.$
		Therefore, there exists $i_0 \neq j_0$ such that $\alpha_{i_0 j_0} \alpha_{j_0 i_0} \geq \dfrac{\sum\limits_{i \neq j}\alpha_{ij} \alpha_{ji}}{N(N-1)} = \dfrac{\mu}{N(N-1)} \geq 0,$ by Theorem \ref{thm5point2} and \eqref{eqation7}. Therefore, we obtain
\[
r_{2}(F,G) \geq \left| \dfrac{\operatorname{tr}(K)}{N} + \sqrt{ \alpha_{i_0 j_0} \alpha_{j_0 i_0} } \right| \geq \dfrac{\operatorname{tr}(K)}{N} + \sqrt{ \dfrac{\mu}{N(N-1)} } = \dfrac{\operatorname{tr}(K)}{N} + \sqrt{ \dfrac{N\,\operatorname{tr}(K^2) - \left( \operatorname{tr}(K) \right)^2 }{N^2(N-1)} }.
\]
Since the right-hand side of the inequality does not depend on the choice of the dual pair \( (F,G) \), it gives a lower bound for \( \tilde{r}_{2} \). On the other hand, suppose $\mu <0.$ Then, by \eqref{eqation7}, for any $K-$dual pair $(F,G),\;$ $\min\limits_{i \neq j} \alpha_{ij}\alpha_{ji} \leq \dfrac{\mu}{N(N-1)}.$  Therefore, there exists $i_1 \neq j_1$ such that $\alpha_{i_1 j_1} \alpha_{j_1 i_1} \leq \dfrac{\sum\limits_{i \neq j}\alpha_{ij} \alpha_{ji}}{N(N-1)} = \dfrac{\mu}{N(N-1)} \leq 0,$ by Theorem \ref{thm5point2} and \eqref{eqation7}. So, 
$$r_{2} (F,G) \geq  \left| \dfrac{tr(K)}{N} + \sqrt{ \alpha_{i_1 j_1}\alpha_{j_1 i_1}} \right| =\sqrt{\left(\dfrac{tr(K)}{N}\right)^2 - \alpha_{i_1 j_1}\alpha_{j_1 i_1}}\geq \sqrt{\left(\dfrac{tr(K)}{N}\right)^2 - \dfrac{\mu}{N(N-1)} }  = \sqrt{\dfrac{(tr(K))^2 -  tr(K^2)}{N(N-1)}}.$$
		
		
		\par 	Assume that there exists a $2$-uniform $K$-dual pair $(F'', G'')$. By applying equation~\eqref{eqn3point4*}, we obtain
\[
r_2(F'', G'') = \left| \dfrac{\operatorname{tr}(K)}{N} + \sqrt{c_{F'', G''}} \right|,
\]
where \( c_{F'', G''} \) denotes the constant arising from the $2$-uniform property of the pair \( (F'', G'') \). This constant also satisfies the identity
\[
\sum_{i \neq j} \alpha''_{ij} \alpha''_{ji} = N(N-1) c_{F'', G''},
\]
which leads to
\[
c_{F'', G''} = \dfrac{\mu}{N(N-1)}.
\]
Based on the sign of \( \mu \), we obtain the following expression for \( \tilde{r}_2 \):
\[
\tilde{r}_2 =
\begin{cases}
\dfrac{\operatorname{tr}(K)}{N} + \sqrt{\dfrac{N\,\operatorname{tr}(K^2) - (\operatorname{tr}(K))^2}{N^2(N-1)}} & \text{if } \operatorname{tr}(K^2) \geq \dfrac{(\operatorname{tr}(K))^2}{N}, \\[2em]
\sqrt{\dfrac{(N-2)(\operatorname{tr}(K))^2 + N\,\operatorname{tr}(K^2)}{N^2(N-1)}} & \text{if } \operatorname{tr}(K^2) < \dfrac{(\operatorname{tr}(K))^2}{N}.
\end{cases}
\]
From this, it can be concluded that \eqref{equation7expression} holds.
\par From the arguments above, we infer that any  $2-$uniform $K-$dual pair $(F,G)$ will satisfy ${r}_{2} (F,G) = \tilde{r}_{2}$ and hence will belong to $\mathcal{R}_{2},$ by Theorem \ref{thm5point2}.	 Conversely, if $(F,G) \in \mathcal{R}_{2},$ then it is a $1-$uniform $K-$dual pair, as it is also in $\mathcal{R}_{1}.$ We shall show that $\alpha_{k\ell}\alpha_{\ell k} = \dfrac{\mu}{N(N-1)},\,\forall\;k \neq \ell.$ Let us first consider the case when  $\mu \geq 0.$ Suppose for some $k \neq \ell,\, \alpha_{k\ell}\alpha_{\ell k} >  \dfrac{\mu}{N(N-1)}.$ Then, using \eqref{eqn3point4*},
		$${r}_{2} (F,G) \geq  \left| \dfrac{tr(K)}{N}+ \sqrt{ \alpha_{k\ell}\alpha_{\ell k}}   \right| > \dfrac{tr(K)}{N} + \sqrt{ \frac{\mu}{N(N-1)}} = \tilde{r}_{2},$$
		which contradicts the fact that $(F,G) \in \mathcal{R}_{2}.$ Therefore, $ \alpha_{k\ell}\alpha_{\ell k} \leq \dfrac{\mu}{N(N-1)},$	for all $k \neq \ell. $ If $\alpha_{k_0 \ell_0}\alpha_{\ell_0 k_0} < \dfrac{\mu}{N(N-1)},$ for some $k_0 \neq \ell_0,$ then $\mu =\displaystyle{\sum_{k \neq \ell}\alpha_{k \ell}\alpha_{\ell k} < \frac{\mu}{N(N-1)} N(N-1)} = \mu,$ thereby proving our claim. A similar argument can be applied to handle the case when \( \mu < 0 \).
	\end{proof}

\end{document}
